% !TeX root = ../kazymyr-thesis.tex

\chapter{Introduction}\label{chapter:introduction}

\section{Motivation}\label{sec:intro-motivation}

Engineering systems maintain several models of the same product side by
side: CAD geometry, a simulation mesh, a written specification. These
models overlap semantically, in the sense that the same underlying fact is
represented in more than one of them. Keeping the representations in
agreement is automated through Consistency Preservation Rules (CPRs),
which take an update made in one model and propagate it into every model
that shares the affected information.

The open question is \emph{when} that propagation is scheduled. Under ACID
semantics the CPR work sits on the critical path: a transaction does not
commit until every consequence of its write has been applied, so the
models are never observably inconsistent, and throughput pays for it.
Under BASE semantics the CPR work is handed to background reconcilers.
Foreground throughput is high, but the models are transiently
inconsistent, and how inconsistent they get is not something the operator
chooses.

In both modes the realised level of inconsistency is an outcome of the
configuration rather than an input to it. ACID pins it to zero, BASE lets
the reconcilers settle it, and neither takes a bound the operator states
in advance. That is the gap this thesis addresses.

\section{Problem statement}\label{sec:intro-problem}

Under pure BASE, the inconsistency level is an emergent property of the
configuration. The only lever available is the number of reconciler
processors, $R$. It is coarse, indirect, and it costs foreground
throughput across the whole system. Under pure ACID, inconsistency is zero
by construction and throughput collapses. In neither case can an operator
ask for a given consistency level and have the system deliver it.

This thesis introduces a HYBRID execution mode, in which the choice
between ACID and BASE is made per transaction, at admission time, by a
component called a \emph{mode oracle}. The adaptive oracle makes that
choice from a threshold on the measured dirty fraction. The claim under
test is that this turns inconsistency from an emergent outcome into a
configurable target.

\section{Formal problem definition}\label{sec:intro-formal}

At each step of a run the system carries some level of inconsistency,
because an update that one model owes another can be accepted now and
applied later. That level is the \emph{dirty fraction} $d_k$, the share of
the resources of the system that are waiting for such an update at step $k$
(Section~\ref{sec:prelim-inconsistency}). A run therefore produces a profile
$\{d_k\}$ over its whole length and not a single number.

That profile is an outcome of the configuration and not an input to it.
Under pure ACID $d_k$ is zero at every step, because nothing is ever left
pending. Under pure BASE it settles wherever the workload and the number of
reconciler processors leave it. Neither mode takes a value for it from the
operator.

The task of this thesis is to build a mechanism that does. Let $c$ be the
inconsistency the operator is willing to carry, stated in advance and in the
same unit as $d_k$. The mechanism has to hold the dirty fraction under $c$
for at least 95\% of the run, which is

\begin{equation}
      p_{95}(d) \;\le\; c + \varepsilon
      \label{eq:problem}
\end{equation}

for a small tolerance $\varepsilon$, where $p_{95}(d)$ is the 95th
percentile of $\{d_k\}$. It has to hold with one setting of its own
parameters rather than with a setting retuned for every configuration of the
system. Throughput is what is paid for that, and the less that is paid the
better.

What the mechanism decides, which is the execution mode of each arriving
transaction, is developed in Section~\ref{sec:approach-hybrid}. How that
decision is made is the subject of Section~\ref{sec:approach-oracles}.

\section{Research questions}\label{sec:intro-rqs}

Three research questions follow from the problem above. The first one asks
whether the hybrid mode can be built at all, the second one what it costs
and whether it is controllable, and the third one whether the feedback
signal is what makes it work.

\begin{description}
      \item[RQ1] How can a multi-model transactional system integrate a
            \emph{hybrid mode} that safely adapts between ACID and BASE semantics
            at runtime, when its subsystems hold overlapping resources?

      \item[RQ2] How does the hybrid mode compare with pure ACID and pure BASE
            execution in throughput and inconsistency, and can the inconsistency it
            permits be configured?

      \item[RQ3] Does choosing the mode from the system's own internal metrics
            outperform a fixed ACID/BASE split on the same
            throughput/inconsistency trade-off?

\end{description}

In the terms of Section~\ref{sec:intro-formal}, RQ1 asks whether such a
mechanism can be built at all, RQ2 asks what it costs and whether $c$
controls the bound in Equation~\ref{eq:problem}, and RQ3 asks whether the
decision has to be taken from the measured inconsistency or whether a fixed
rule suffices.

\section{Overview of the approach}\label{sec:intro-approach-overview}

The approach has three parts. First, a HYBRID execution mode extends the
base simulator so that each root transaction is tagged \textsc{acid} or
\textsc{base} at admission and executes accordingly, with both paths
sharing the same resources, processors and CPR machinery. Second, because
the two paths now coexist, an arbiter defines what happens when they meet
on the same resource: an ACID transaction that needs a dirty resource
waits for reconciliation to clear it, and a BASE transaction that would
dirty a locked resource defers the repair and commits anyway, so
nothing on the BASE side ever blocks. This single rule makes the mixed
state well defined, and because the wait edge it adds points one way only,
it introduces no new deadlock. Third, the tagging decision is delegated to
a mode oracle. The adaptive oracle switches to ACID when the dirty
fraction, read live at the moment of the decision, reaches the configured
threshold $c$, and back to BASE when it falls to a lower recovery
threshold $r$; the interval $[r, c]$ is a hysteresis band, and no
smoothing filter is applied to the signal, for reasons the evaluation
establishes (Section~\ref{sec:approach-nofilter}). The fixed-share oracle,
which holds the ACID share of transactions at a configured $p_{\mathrm{acid}}$ and reads
nothing, is the baseline that isolates what the feedback contributes.

The evaluation follows the same logic. Five parameter sweeps,
640 simulation runs in total, establish that the
mixed state is reached and resolved constantly rather than never, where
the oracle engages, what the mixture costs, whether $c$ controls the
realised inconsistency, whether a fixed oracle can do the same job, and
over which range of workloads the answer holds.

\section{Contributions}\label{sec:intro-contributions}

The work is built on the multi-model simulator of Zacouris and
Acosta~\cite{zacouris2025m2ts,zacouris2026mxs}, which runs ACID or BASE but
never both. What this thesis adds to it is the following.

\begin{bullets}
  \item[A hybrid execution mode.] The execution mode becomes a tag carried
  by each root transaction instead of a setting fixed for the whole run.
  Both paths share the same resources, processors and CPR machinery
  (Section~\ref{sec:approach-hybrid}).

  \item[A shared-resource arbiter.] Once both paths run together, a
  resource can be ACID-locked and BASE-dirty at the same time, which
  neither pure mode ever produces. The arbiter forbids that state. An ACID
  transaction waits for the repair, and a BASE transaction defers its own
  repair and commits (Section~\ref{sec:approach-arbiter}).

  \item[A safety argument.] The wait edge the arbiter adds between the two
  modes points one way only, so the mixed state adds no deadlock that pure
  ACID did not already admit (Section~\ref{sec:approach-deadlock}).

  \item[Two mode oracles.] The adaptive oracle reads the measured dirty
  fraction and turns ACID on at a configured threshold, with a hysteresis
  band below it. The fixed-share oracle reads nothing and holds the ACID
  share at a configured value. It is the baseline that isolates what the
  feedback signal contributes (Section~\ref{sec:approach-oracles}).

  \item[Metrics for mixed runs.] A hybrid run is measured by its mode
  switches, its ACID share by time and by transaction count, its cross-mode
  blocking and the duration of its ACID visits. These separate a real
  mixture of the two modes from a run that never left one of them.

  \item[An empirical evaluation.] Five parameter sweeps over the system and
  over the workload, 640 simulation runs in total, every configuration
  replicated under five seeds (Chapter~\ref{chapter:experiments}).
\end{bullets}

\section{Structure of the thesis}\label{sec:intro-structure}

Chapter~\ref{chapter:preliminaries} introduces the concepts the thesis
relies on: transactions and two-phase locking, BASE and reconciliation,
multi-model systems and CPRs, the base simulator this work extends, and
the dirty set as the control signal the oracle reads.
Chapter~\ref{chapter:related-work} positions the work against the
literature, nearest first: the simulator this thesis extends, the systems
that already choose an execution mode per transaction at runtime, and then
the work on combining ACID and BASE, on switching consistency at runtime
and on bounded inconsistency.
Chapter~\ref{chapter:approach} describes the approach: the hybrid mode,
the correctness problem it creates, the arbiter that resolves it, and the
two oracles. Chapter~\ref{chapter:experiments} gives the experimental
settings, validates the arbiter against the run data, reports the five
sweeps, and closes with two analyses that apply to all of them: where the
oracle actually leaves ACID as against where it was configured to, and why
the smoothing filter was removed.
Chapter~\ref{chapter:conclusion} answers each research question, records
the lessons learned, and lists the open questions.
